\begin{table}[H]\centering
\caption{\textbf{Internal consistency of the cesarean-timing indicator across Robson groups}}
\label{tab:robson_validation}
\small
\small\setlength{\tabcolsep}{4pt}\resizebox{\ifdim\width>\linewidth \linewidth\else\width\fi}{!}{%
\begin{tabular}{lrr}
\toprule
Robson group & Cesareans & Share coded prelabor (\%) \\
\midrule
01 & 1,499,945 & 0.0 \\
02 & 2,211,444 & 84.7 \\
03 &   694,312 & 0.0 \\
04 &   922,035 & 85.0 \\
05 & 3,780,923 & 56.5 \\
06 &   261,292 & 53.0 \\
07 &   340,458 & 52.0 \\
08 &   376,493 & 56.7 \\
09 &    45,713 & 49.8 \\
10 & 1,058,662 & 53.8 \\
\midrule
\multicolumn{3}{p{0.9\linewidth}}{\footnotesize Group 1 (spontaneous labor) has a 0.0\% prelabor share, as it must; its for-profit weekend dip of -6.5 pp is therefore entirely in-labor (in-labor component -5.5 pp). A formal test cannot distinguish the group-1 dip from the group-10 (preterm) dip ($p=0.18$), so the preterm placebo is weak.} \\
\bottomrule
\end{tabular}
}
\\[2pt]\footnotesize\textit{Notes:} SINASC 2010--2024, cesareans with a valid before/during-labor code. Robson 1 = nulliparous, term, singleton, cephalic, spontaneous labor; Robson 2 = same but induced or prelabor cesarean.
\end{table}
